\section{Open questions}

The following five questions are genuinely open --- neither the paper nor
this replication answers them. Each includes a concrete free-endpoint
next-step probe.

\begin{enumerate}

\item \textbf{Problem-adaptive bandwidth without a full grid sweep.}
Does a problem-adaptive $\lambda$ (chosen per training set via kernel-target
alignment or a Silverman-style scale rule) outperform the paper's single
grid-scanned $\lambda^*$, and can it be chosen without a full grid sweep?
\emph{Basis:} the paper picks $\lambda$ by coarse validation grid; no
analytic prescription or adaptive-$\lambda$ algorithm is proposed, and the
linkage to classical bandwidth-selection theory (Silverman, kernel alignment,
MMD) is left implicit.
\emph{Next step:} on 3 datasets (\texttt{make\_moons}, digits-PCA-4d, wine),
compare (a) grid-search $\lambda^*$, (b) kernel-target-alignment maximizer
$\lambda_{\text{KTA}}$, (c) Silverman analog $\lambda_{\text{Sil}} = c \cdot n^{-1/(d+4)} \cdot \mathrm{std}(X)$.
Score by test accuracy and wall-clock cost. Free: local PennyLane, $\sim$1\,hr.

\item \textbf{Joint quantum-kernel-alignment + bandwidth optimization.}
Does integrating bandwidth optimization with quantum-kernel alignment
(Glick et al.\ 2022, Hubregtsen et al.\ 2022) yield a jointly-optimal
$(\theta, \lambda)$, and does the paper's optimum sit inside the joint
sweet-spot?
\emph{Basis:} the paper treats the feature map as fixed. Modern QKA
jointly optimizes ansatz $\theta$ AND embedding. Interaction is unexplored.
\emph{Next step:} on \texttt{make\_moons} + a 4-qubit trainable ansatz
(\texttt{HardwareEfficient(reps=2)}), grid-sweep $(\theta_{\text{aligned}}(\lambda), \lambda)$
in 2D; compute alignment landscape $A(\theta, \lambda)$; check whether
$\operatorname{argmax}_\lambda$ Shaydulin-Wild (with $\theta$ frozen at random
init) coincides with the joint $\operatorname{argmax}$. Free: $\sim$20\,min.

\item \textbf{Chemistry / hard-vision generalization.}
Does the bandwidth sweet-spot survive on genuinely-hard tasks (chemistry PES,
Fashion-MNIST-full, CIFAR-10-PCA), or does it require datasets that are
already easy for classical kernels?
\emph{Basis:} \texttt{make\_moons} is 2D and near-linearly-separable; the
paper's chosen datasets (Fashion-MNIST-PCA, KMNIST-PCA, PLASTiCC) are all
PCA-preprocessed. On genuinely-hard tasks (SMILES property regression,
bandgap prediction) the accuracy curve may be flatter, shifted, or absent.
\emph{Next step:} run Shaydulin-Wild sweep on QM9-subset ($\sim$200 molecules,
target = HOMO-LUMO gap, features = Coulomb-matrix eigenvalues PCA'd to 4d)
and separately on PLASTiCC subset if PCA loadings are released. Look for
whether $\lambda^*$ still exists and moves systematically with dataset
intrinsic dimension. Free: local PennyLane + QM9 subset from PyG/MoleculeNet.

\item \textbf{Noise robustness of the sweet-spot.}
How does $\lambda^*$ shift under realistic shot noise, depolarizing noise, or
coherent gate errors --- is the noiseless-optimum still optimal at
1000 shots + typical superconducting T1/T2, or does noise \emph{shift}
$\lambda^*$?
\emph{Basis:} the paper's sims are exact-statevector, no shots, no gate noise.
Real quantum-kernel evaluation is sample-complexity-bound ($O(1/\epsilon^2)$
shots per matrix element) and noise-corrupted (depolarizing flattens fidelity
toward $1/2^n$, which itself concentrates the kernel). Shot noise adds a floor
to $\mathrm{Var}[K_{ij}]$ that interacts with bandwidth.
\emph{Next step:} rerun sweep with \texttt{default.mixed}, add depolarizing
$p \in \{0, 0.001, 0.01, 0.05\}$ per gate, and shots $\in \{100, 1000, 10000, \infty\}$.
Plot $\lambda^*$ vs.\ $(p, \text{shots})$. \emph{A priori} prediction: noise
should shift $\lambda^*$ smaller (larger bandwidth) because noise-induced
kernel flattening compensates for bandwidth-induced concentration. Free:
$\sim$2\,hr PennyLane.

\item \textbf{Analytic $\lambda^*$ from expressivity bounds.}
Is the sweet-spot related to modern QML expressivity bounds
(K\"ubler et al., Thanasilp et al.\ exponential concentration, Larocca et al.\
barren-plateau bounds), and can $\lambda^*$ be derived \emph{analytically}
from the feature-map's Lipschitz constant + data covariance instead of
grid-searched?
\emph{Basis:} the paper cites the concentration no-go as motivation but does
not derive a closed-form $\lambda^*(n, \text{data}, \text{feature-map})$.
Subsequent work gives concentration bounds that scale with an effective
Lipschitz constant of the embedding. In principle these should \emph{predict}
the sweet-spot.
\emph{Next step:} compute $L(\lambda)$ for the IQP embedding analytically
(bound via $\|\mathrm{d}\phi/\mathrm{d}x\| = O(\lambda)$ for depth-$r$ IQP);
use Thanasilp et al.\ concentration bound
$\mathrm{Var}\,K = O(e^{-c n L^2 \lambda^2})$ to predict
$\lambda_{\text{predicted}}^*$ by setting the bound to a fixed threshold
(e.g.\ $\mathrm{Var}\,K = 0.1$ of its max). Compare
$\lambda_{\text{predicted}}^*$ to empirical $\operatorname{argmax}_\lambda$
across our sweep and across the paper's Fig.\ 2 panels. Free: half day
pen-and-paper + PennyLane cross-check.

\end{enumerate}
